% !PRINTED 460
\weberpage{460}
%<<CONT

\[
\begin{aligned}
D &= -8, & \omega &= \sqrt{-2}, \\
D &= -7, & \omega &= \frac{-3+\sqrt{-7}}{2}, \\
D &= -4, & \omega &= i.
\end{aligned}
\]

Der Wert $\gamma_2(i)$ ist aber (nach \S\,117) bekannt, nämlich
$=12$, und daher muß die linke Seite von (14) durch $u-12$ teilbar
sein. Da dieser Faktor bekannt ist, findet man leicht die übrigen:

\[
u^4-2u^3-495u^2+2^4\cdot 3^3\cdot 5^3 = (u-12)(u-20)(u+15)^2.
\]

Da $\gamma_2(\omega)$ ebenso wie $j(\omega)$ für ein rein imaginäres
$\omega$ einen positiven Wert hat, so muß der Faktor $u-20$ dem
Werte $D = -8$ entsprechen, und wir erhalten:

\weq{16}{\gamma_2(i) = 12,}

\weq{17}{\gamma_2(\sqrt{-2}) = 20,}

\weq{18}{\gamma_2\left(\frac{-3+\sqrt{-7}}{2}\right) = -15.}

Es existieren außer $-7$ noch fünf ungerade durch 3 nicht teilbare
negative Diskriminanten:

\weq{19}{-11, \quad -19, \quad -43, \quad -67, \quad -163^{1)},}

für die es nur \gesperrt{eine} Formenklasse gibt. Für diese ist also nach
unserem Satze:

\[
-\gamma_2\left(\frac{-3+\sqrt{-m}}{2}\right)
\]

\gesperrt{eine ganze rationale Zahl} $Z$.

Um diese rationalen Zahlen zu finden, wollen wir Grenzen aufsuchen,
zwischen denen sie liegen müssen, die um weniger als eine Einheit
voneinander verschieden sind.

Nach \S\,54, (5), (8) ist:

\[
\gamma_2(\omega) = f(\omega)^{16}-\frac{16}{f(\omega)^8} = f_1(\omega)^{16}+\frac{16}{f_1(\omega)^8} = f_2(\omega)^{16}+\frac{16}{f_2(\omega)^8},
\]

ferner nach \S\,34, (19):

\[
f_2\left(\frac{-3+\omega}{2}\right) f(\omega) = e^{-\frac{\pi i}{8}}\sqrt{2},
\]

\bigskip
\noindent\rule{4cm}{0.4pt}

\footnotesize
1) Daß nicht mehr Diskriminanten dieser Art existieren, kann bis
jetzt nur daraus geschlossen werden, daß, soweit man die Berechnung
der Klassenzahlen fortgesetzt hat, weitere Zahlen dieser Art nicht
gefunden sind [vgl. die Tafel der Klassenzahlen von
\gesperrt{Gauss} (Werke, Bd.\,II, S.\,450)].
\normalsize

%CONT>>
